\section{Study Design}
\label{sec:study}

\subsection{Research Questions}
\label{sec:study:rqs}

\begin{description}
  \item[RQ1] \emph{Integrated workspace value (baseline).}
    Does using an integrated workspace (without an agent) reduce cognitive load
    and time-on-task relative to reported experience with fragmented tooling?

  \item[RQ2] \emph{Effect of context breadth on agent effectiveness.}
    Does increasing the agent's context breadth (\CI $\to$ \CII $\to$ \CIII)
    improve task performance and reduce cognitive load?

  \item[RQ3] \emph{Task-type moderation.}
    Does the effect of context breadth vary across task types (comprehension,
    single-surface, cross-surface)?

  \item[RQ4] \emph{Scaling of benefit with surface integration.}
    Do participants in higher-context conditions show greater cross-surface
    integration behavior (measured by telemetry) and higher integration
    perception scores?
\end{description}

\subsection{Hypotheses}
\label{sec:study:hypotheses}

Based on the research questions and prior work on AI-assisted workflows, we
state four directional hypotheses:

\begin{description}
  \item[H1~(Context breadth and speed).]
    Broader agent context reduces time-on-task: participants in \CIII\ will
    complete agent tasks faster than those in \CII, who will in turn be faster
    than those in \CI.

  \item[H2~(Task-type moderation).]
    The effect of context breadth is moderated by task type. The advantage of
    \CIII\ over \CI\ will be negligible for comprehension and single-surface
    tasks (where cross-surface data is irrelevant) but substantial for
    cross-surface tasks that require coordinating board, chat, and GitHub.

  \item[H3~(Cognitive load).]
    Broader context reduces perceived cognitive load (\NASATLX): \CIII\ $<$
    \CII\ $<$ \CI, driven primarily by effort and mental-demand subscales.

  \item[H4~(Perception--behavior gap).]
    Trust and integration-perception scores will not fully track objective task
    performance: participants may report high trust or high integration
    perception even in conditions where telemetry shows limited cross-surface
    coordination, and vice versa.
\end{description}

\subsection{Participants}
\label{sec:study:participants}

% TODO: Fill in once recruitment is complete.

We target $N$~=~20--30 participants drawn from two complementary populations:
\emph{university CS students} (undergraduate and graduate, recruited via
course mailing lists) and \emph{professional software developers} ($\geq$1 year
of industry experience, recruited via professional networks and online
communities). Recruiting both groups allows us to examine whether the benefit of
context breadth varies with prior experience. Eligibility further required
self-reported familiarity with at least one task-management tool (e.g., Jira,
Linear, GitHub Projects) and one team-chat tool (e.g., Slack, Teams, Discord).
The study was conducted remotely. Participants were compensated
\emph{[amount TBD]}.

\subsection{Conditions}
\label{sec:study:conditions}

Participants were randomly assigned, with balanced block randomization (block
size 3), to one of three between-subjects conditions:

\begin{description}
  \item[\CI~(Chat only, advisory).] The agent observes only team chat history.
    No access to board state or GitHub. Advisory-only: no tools offered to the
    model, so it cannot modify any surface.

  \item[\CII~(Chat + Board, action-capable).] The agent observes chat and the
    full board card state. Available tools: \texttt{create\_card},
    \texttt{move\_card}, \texttt{update\_card}, and \texttt{post\_message}
    (drafts a chat message for the participant to review before sending).

  \item[\CIII~(Chat + Board + GitHub, full action).] The agent observes all
    three surfaces. All \CII\ tools plus \texttt{link\_card\_to\_github},
    enabling it to connect board cards to GitHub issues or pull requests.
\end{description}

\subsection{Tasks}
\label{sec:study:tasks}

The study uses a hybrid design. Each participant completes 8 tasks in fixed
order: two manual \emph{baseline tasks} (C0, no agent available) establishing a
per-person speed and workload anchor, followed by six \emph{agent-assisted
tasks} in their assigned condition. Baseline and agent tasks are
content-disjoint to prevent learning transfer. The agent phase uses the
\texttt{gpt-4o-mini-2024-07-18} model at temperature 0.2 with versioned
condition-specific system prompts pinned for reproducibility.

Task types span both phases:

\begin{description}
  \item[Comprehension.] Synthesize information from one or more surfaces to
    form a judgment (e.g., ``what does the team need most right now?'').
  \item[Single-surface.] Can be completed by acting on a single surface, even
    if cross-surface context helps.
  \item[Cross-surface.] Explicitly requires coordinated action across two or
    more surfaces.
\end{description}

\begin{table}[h]
  \caption{Task design summary. B = Baseline (no agent, C0), T = Agent task.
    Surfaces: Ch = Chat, Bo = Board, GH = GitHub.}
  \label{tab:tasks}
  \begin{tabular}{lllll}
    \toprule
    ID & Phase & Type & Surfaces & Summary \\
    \midrule
    B1 & Baseline & Single-surface & Ch, Bo & Find blocked card from chat; move to In Progress \\
    B2 & Baseline & Cross-surface  & Ch, Bo & Check card status; post reply in chat \\
    \midrule
    T1 & Agent & Comprehension  & Ch, Bo     & Determine next priority; update board \\
    T2 & Agent & Comprehension  & Ch, Bo, GH & Catch-up summary; post in chat \\
    T3 & Agent & Single-surface & Ch         & Draft and post status update in chat \\
    T4 & Agent & Single-surface & Bo         & Set priorities on three backlog cards \\
    T5 & Agent & Cross-surface  & Ch, Bo     & Find bug card from chat; move to In Progress \\
    T6 & Agent & Cross-surface  & Ch, Bo, GH & Link PR to card; announce in chat \\
    \bottomrule
  \end{tabular}
\end{table}

\subsection{Procedure}
\label{sec:study:procedure}

% TODO: Fill in IRB protocol details, session timing.

Each session proceeded as follows:
\begin{enumerate}
  \item Informed consent and eligibility screening.
  \item Orientation to \Weld (un-timed, guided).
  \item \textbf{Baseline phase} (2 tasks, no agent, condition C0): tasks
    presented sequentially; participants self-paced, rated each task before
    advancing.
  \item Phase-transition interstitial: agent panel revealed; \CII/\CIII\
    participants received a brief explanation of available agent actions.
  \item \textbf{Agent phase} (6 tasks): same structure with the agent
    available under the participant's assigned condition (\CI, \CII, or
    \CIII).
  \item Exit survey (\NASATLX, \SUS, integration-perception, AI trust scale,
    open-ended comments, demographics).
  \item Debrief.
\end{enumerate}

Sessions were conducted remotely. Screen recording was \emph{not} collected;
all behavioral data came from in-product telemetry.

\subsection{Measures}
\label{sec:study:measures}

\paragraph{Task performance.}
Task completion (binary, judged from \texttt{task\_final\_state} telemetry) and
time-on-task (seconds between \texttt{task\_shown} and \texttt{task\_completed}
or \texttt{task\_skipped}).

\paragraph{Cognitive load.}
\NASATLX~\cite{hart1988tlx}, six subscales (mental demand, physical demand,
temporal demand, performance, effort, frustration), administered once per phase
in the exit survey.

\paragraph{Usability.}
\SUS~\cite{brooke1996sus}, 10 items, 5-point Likert, yielding a 0--100 score.

\paragraph{Integration perception.}
A five-item scale measuring whether participants perceived workspace surfaces as
connected (e.g., ``I could easily see how chat discussions connected to board
tasks'').

\paragraph{AI trust and effectiveness.}
An eight-item scale (adapted from~\cite{madsen2000measuring}) measuring perceived
agent accuracy, helpfulness, and appropriate scope.

\paragraph{Behavioral telemetry.}
Cross-surface reference events (\texttt{card\_referenced\_from\_chat},
\texttt{pr\_linked\_to\_card}, etc.) as an objective measure of integration
behavior; agent turn counts and tool-call distributions as agent-use indicators.
